\section*{Hoofdstuk \ref{ch:b3:quantum-statistics} --- Kwantumstatistiek}

\begin{solution}{exo:b3:quantum-statistics:1}
(a) $\qty{7.1}{eV}$, $\qty{8.2e4}{K}$. (b) $\qty{3.2}{eV}$,
$\qty{3.7e4}{K}$. (c) Met $m = 3u$: $E_{\text{F}} \approx
\qty{4e-4}{eV}$, $T_{\text{F}} \approx \qty{5}{K}$ (de waarde voor een
ideaal gas; de wisselwerkingen verlagen haar). (d) Elektronen in metalen:
bij elke aardse temperatuur ontaard; $^3$He: pas onder enkele kelvin; $T = \qty{300}{K}$ raakt de fermischaal van een metaal nooit.
\end{solution}

\begin{solution}{exo:b3:quantum-statistics:2}
(a) $\int_0^{E_{\text{F}}}\epsilon\,g\,\dd\epsilon/\int g\,\dd
\epsilon = \tfrac35 E_{\text{F}}$ uit $g \propto \sqrt\epsilon$.
(b) $v_{\text{F}} = \sqrt{2E_{\text{F}}/m_{\text{e}}} =
\qty{1.6e6}{m/s}$; effectief $\sqrt{3/5}\,v_{\text{F}} \approx
\qty{1.2e6}{m/s}$. (c) $\tfrac32 k_{\text{B}}T = \tfrac35
E_{\text{F}}$:
$T \approx \qty{3.3e4}{K}$. (d) Er is nergens heen te gaan: elke lagere
toestand is bezet, dus kan er geen foton worden uitgezonden en kan geen
botsing ze afremmen --- snelheid zonder warmte, door uitsluiting.
\end{solution}

\begin{solution}{exo:b3:quantum-statistics:3}
(a) $P = 0.4 \times \num{8.5e28} \times \qty{7.1}{eV} =
\qty{3.9e10}{Pa} \approx \num{4e5}$ atmosfeer. (b) De coulombaantrekking
van het ionenrooster: het metaal is de wapenstilstand tussen beide
(\cref{exo:b3:quantum-statistics:11}). (c) Samendrukken tilt
$E_{\text{F}}$ steil op ($n^{2/3}$): metalen bieden weerstand. (d) $K \sim
\tfrac53 P \approx \qty{65}{GPa}$ tegen de gemeten \qty{140}{GPa}: de
helft van de stijfheid van koper is zuivere uitsluiting.
\end{solution}

\begin{solution}{exo:b3:quantum-statistics:4}
(a) $T_{\text{c}} \approx \qty{0.4}{\micro K}$ --- het condensaat van JILA
verscheen bij \qty{170}{nK}: dezelfde natuurkunde, met een iets lagere
piekdichtheid. (b) $T_{\text{c}} \propto 1/m$: $\times87 \approx
\qty{35}{\micro K}$. (c) Bij hoge dichtheid en zulke kou \emph{stolt} elk
gas behalve spingepolariseerde waterstof; driedeeltjesbotsingen maken
moleculen. De truc is een ijl, metastabiel gas dat koud genoeg is opdat
$\lambda_T$ de verdunning goedmaakt. (d)
$\lambda_{T_{\text{c}}} = h/\sqrt{2\pi mk_{\text{B}}T_{\text{c}}}
\approx \qty{0.3}{\micro m}$: $n\lambda^3 = \num{e20} \times
\num{2.7e-20} \approx 2.7$, zoals de regel vereist.
\end{solution}

\begin{solution}{exo:b3:quantum-statistics:5}
(a) $\gamma T = \beta T^3$ bij $T^2 = \gamma/\beta$: voor koper is
$T \approx \qty{3}{K}$. (b) $C_{\text{el}}/C \approx
(\pi^2/2)(300/\num{8.2e4})/3 \approx 0.6\%$. (c) Bij \qty{0.1}{K} is de
roosterterm als $T^3$ ingestort: de elektronen dragen $\sim10^3$ keer
meer --- de calorimetrie wordt elektronspectroscopie. (d) $C/T = \gamma + \beta T^2$ is een rechte in $T^2$: het snijpunt $\gamma$ meet de
\omterm{prop:b3:schrodinger-three-dimensions:counting}{toestandsdichtheid} op het ferminiveau en de helling $\beta$ de
debyetemperatuur --- twee getallen per metaal uit één grafiek.
\end{solution}

\begin{solution}{exo:b3:quantum-statistics:6}
(a) $N_{\text{e}} = M/2m_{\text{n}} = \num{6e56}$; $n =
\num{5.5e35}\,\unit{m^{-3}}$; $E_{\text{F}} \approx
\qty{0.24}{MeV}$. (b)
De helft van de \omterm{def:b3:relativistic-dynamics:momentum-energy}{rustenergie} van het elektron (op de rand van
relativistisch), en toch driehonderd keer $k_{\text{B}}T$
(\qty{0.9}{keV}): tegelijk witgloeiend en kwantumkoud. (c)
$P_{\text{deg}} \approx \qty{9e21}{Pa}$ tegen $GM^2/R^4 \sim
\num{e22}\,\unit{Pa}$: dezelfde orde --- de ster is precies dat
evenwicht. (d) Meer massa: de kant van de zwaartekracht groeit als $M^2$,
en de ster moet krimpen ($R \propto M^{-1/3}$), wat $E_{\text{F}}$ naar de
relativistische verzachting drijft.
\end{solution}

\begin{solution}{exo:b3:quantum-statistics:7}
(a) $\rho \approx \qty{5e17}{kg/m^3}$: ongeveer tweemaal de
kerndichtheid --- een atoomkern zo groot als een stad. (b) $n_{\text{n}}
\approx \qty{3e44}{m^{-3}}$: $E_{\text{F}} \approx \qty{90}{MeV}$, een
tiende van $m_{\text{n}}c^2$: de relativiteitstheorie klopt aan. (c) De
\omterm{thm:b3:quantum-statistics:fermi}{fermi-energie} van de elektronen steeg voorbij de prijs van
\qty{0.78}{MeV} van $\text{e} + \text{p} \to \text{n} + \nu$: de vangst
werd lonend, en de ster ruilde haar elektronen en protonen voor
neutronen. (d) $\qty{5}{cm^3} \times \qty{5e17}{kg/m^3} =
\qty{2.5e12}{kg}$: een theelepel die een bergketen overtreft.
\end{solution}

\begin{solution}{exo:b3:quantum-statistics:8}
(a) Een holte afkoelen \emph{verwijdert} eenvoudigweg fotonen (de wanden
slikken ze op): zonder een aantal om te behouden blijft $\mu$ op nul
vastgepind en hoopt er zich nooit een overschot in de grondmodus op --- het
gas dooft in plaats van te condenseren. (b) Het aantal atomen blijft op
elke laboratoriumtijdschaal behouden. (c) De kleurstof absorbeert en
zendt de fotonen herhaaldelijk uit \emph{zonder ze te verliezen} op de
thermalisatietijdschaal: een in de praktijk behouden aantal fotonen (en een
effectieve massa uit de holte) --- en de condensatie verschijnt dan ook.
(d) Fononen: vrijelijk geschapen en vernietigd, $\mu = 0$: geen
condensatie --- het ``fononengas'' van een kristal vriest gewoon uit.
\end{solution}

\begin{solution}{exo:b3:quantum-statistics:9}
(a) $T_{\text{c}}^{\text{ideaal}} \approx \qty{3.1}{K}$ tegen de gemeten
$T_\lambda = \qty{2.17}{K}$. (b) Vloeibaar helium is dicht en sterk
wisselwerkend: ``ideaal'' laat juist de wisselwerkingen vallen die de
overgang verschuiven (en van vorm veranderen). (c) De atomen van $^3$He
zijn fermionen: geen condensatie zonder paarvorming; aantrekkende
wisselwerkingen binden cooperachtige paren pas onder \qty{2.6}{mK}, en de
gepaarde vloeistof stroomt dan superieur. (d) Negen fermionen: $^6$Li is
een fermion (terwijl $^7$Li een boson is) --- met het paar lithiumisotopen
kan één laboratorium een fermizee en een condensaat in dezelfde oven
bestuderen.
\end{solution}

\begin{solution}{exo:b3:quantum-statistics:10}
(a) Vul $g \propto \sqrt\epsilon$ in en schaal $\beta$ eruit. (b) De
waarde van de integraal geeft $N_{\text{exc}} =
2.612\,V/\lambda_T^3$;
$N_{\text{exc}} = N$ stellen definieert $T_{\text{c}}$. (c) Onder
$T_{\text{c}}$ is $N_{\text{exc}}(T) =
N(T/T_{\text{c}})^{3/2}$; de rest
is $N_0$. (d) In twee dimensies is $g(\epsilon)$ constant en divergeert
$\int\dd\epsilon/(\eu^{
\beta\epsilon} - 1)$ aan de onderkant: de
aangeslagen toestanden kunnen iedereen altijd opnemen --- geen
macroscopische bezetting van de grondtoestand bij enige $T > 0$.
\end{solution}

\begin{solution}{exo:b3:quantum-statistics:11}
(a) $\dd\epsilon/\dd n = \tfrac23 An^{-1/3} - \tfrac13 Bn^{-2/3} =
0$ bij $n^{1/3} = B/2A$: een echt minimum. (b) Met $A \sim
\hbar^2/m_{\text{e}}$ en $B \sim e^2/4\pi\varepsilon_0$ is de
evenwichtsafstand $n^{-1/3} \sim \hbar^2 4\pi\varepsilon_0/
m_{\text{e}}e^2 = a_0$: metalen zijn dicht omdat de \omterm{thm:b3:hydrogen-atom:levels}{bohrstraal} dat zegt.
(c) $\epsilon_{\min} \sim -B^2/A \sim -E_{\text{I}}$ --- cohesie op de
schaal van elektronvolts, zoals gemeten. (d) De \omterm{thm:b3:quantum-statistics:fermi}{ontaardingsdruk} is de
afstotende helft van de wapenstilstand: zonder haar zou de coulombterm
het rooster tot een punt verpletteren --- het volume van de materie is
dat van Pauli.
\end{solution}

\begin{solution}{exo:b3:quantum-statistics:12}
(a) Actieve elektronen: $g(E_{\text{F}})k_{\text{B}}T$; energie elk: $\sim k_{\text{B}}T$; differentieer $E_{\text{extra}} \sim
g k_{\text{B}}^2T^2$. (b) $g(E_{\text{F}}) = 3N/2E_{\text{F}}$ geeft $C \sim 3Nk_{\text{B}}T/T_{\text{F}}$ --- de exacte $\pi^2/2$ vervangt de 3.
(c) Omklapbare spins $\propto T$; elk draagt een uitlijning $\propto \mu^2B/k_{\text{B}}T$ bij: de $T$'s heffen elkaar op --- de
susceptibiliteit van Pauli is vlak waar die van Curie
(\cref{exo:b3:canonical-ensemble:11}) als $1/T$ daalt: een diagnose van
ontaarding. (d) Ontaarde fermionen reageren alleen aan hun oppervlak:
vermenigvuldig elk klassiek antwoord met $\sim
T/T_{\text{F}}$, of beperk
het tot de fermischil.
\end{solution}

\begin{solution}{pb:b3:quantum-statistics:1}
\textbf{1.} De energie die vrijkomt bij het samenstellen van de ster uit
verspreide materie --- negatief, en dieper naarmate zij compacter is.
\textbf{2.} $E_{\text{F}} \sim (\hbar^2/m_{\text{e}})
(N_{\text{e}}/R^3)^{2/3}$.
\textbf{3.} $E_{\text{k}} \sim N_{\text{e}}E_{\text{F}} \sim
\hbar^2N_{\text{e}}^{5/3}/m_{\text{e}}R^2$.
\textbf{4.} $A/R^2 - B/R$ minimaliseren geeft $R = 2A/B$: $R \sim
\hbar^2N_{\text{e}}^{-1/3}/(Gm_{\text{e}}m_{\text{n}}^2\mu_e^2)$.
\textbf{5.} $N_{\text{e}} \propto M$: $R \propto M^{-1/3}$ --- massa
toevoegen laat de ster \emph{krimpen}, zodat de dichtheid en
$E_{\text{F}}$ onbegrensd groeien: de niet-relativistische steun leeft op
krediet.
\textbf{6.} $N_{\text{e}} = \num{6e56}$: $R \sim \qty{2000}{km}$ --- de
juiste orde naast de \qty{5800}{km} van Sirius B (onze schaling liet
factoren van enkele eenheden vallen).
\textbf{7.} $E_{\text{F}} \sim m_{\text{e}}c^2$ bij $n \sim
(m_{\text{e}}c/\hbar)^3 \approx \qty{2e37}{m^{-3}}$, dus $\rho
\sim \num{e10}$--$\num{e11}\,\unit{kg/m^3}$: het regime van de zwaarste
dwergen.
\textbf{8.} Elk elektron: $\epsilon \sim \hbar c\,n^{1/3}$; in totaal:
$\hbar cN_{\text{e}}^{4/3}/R$.
\textbf{9.} Met beide termen $\propto 1/R$ valt de straal weg: de
stabiliteit wordt beslist door de vraag of de kinetische coëfficiënt de
gravitationele overtreft --- een zuiver massacriterium.
\textbf{10.} Gelijkstellen: $N_{\text{e,crit}} \sim (\hbar c/
Gm_{\text{n}}^2\mu_e^2)^{3/2}$, en $M_{\text{Ch}} =
\mu_em_{\text{n}}N_{\text{e,crit}}$: de aangegeven formule.
\textbf{11.} $\hbar c/Gm_{\text{n}}^2 = \num{1.7e38}$;
$M_{\text{Ch}} \sim (\num{1.7e38})^{3/2}m_{\text{n}}/4 \approx
\qty{9e29}{kg} \approx 0.5\,M_\odot$ --- de exacte berekening scherpt de
orde van grootte aan tot $1.4\,M_\odot$.
\textbf{12.} $\hbar$ (de druk is kwantummechanisch), $c$ (haar
relativistische verzachting) en $G$ (de tegenstander): een getal op
sterschaal, uit drie microscopische constanten opgebouwd --- de
kwantummechanica die wetten stelt voor voorwerpen van $10^{57}$ deeltjes.
\textbf{13.} De ontaarding van de neutronen, bij een straal die met
$\sim m_{\text{e}}/m_{\text{n}}$ kleiner is (maal de boekhouding van
$\mu_e$): kilometers --- de neutronenster.
\textbf{14.} Dezelfde drie constanten met overal $m_{\text{n}}$ geven
dezelfde massaschaal: neutronensterren houden op nabij $2\,M_\odot$;
daarvoorbij wint geen enkele drukwet --- een zwart gat.
\textbf{15.} $GM^2/R \approx \qty{5e46}{J}$: enkele honderden malen de
volledige opbrengst van de zon over tien miljard jaar, in seconden
vrijgemaakt.
\textbf{16.} Supernova 1987A: twee dozijn neutrino's, uren voordat het
licht opflakkerde aangekomen en precies de voorspelde energie dragend ---
de ineenstorting rechtstreeks gehoord.
\textbf{17.} Kerninstorting: de ijzeren kern van een zware ster implodeert
voorbij elke drukbodem; de mantel wordt weggeblazen. Type Ia: een \omterm{ex:b3:quantum-statistics:dwarf}{witte
dwerg} die tot aan de drempel is gevoed ontsteekt zijn koolstof in een
thermonucleaire flits die de hele ster ontbindt.
\textbf{18.} Dezelfde ontstekingsmassa, dezelfde brandstof, dezelfde
energieafgifte: vrijwel identieke lichtkrommen --- geijkte bommen.
\textbf{19.} Kennen we de absolute helderheid, dan geeft de waargenomen
helderheid de afstand volgens de omgekeerde kwadratenwet; zonder
``standaard'' stort de ladder in.
\textbf{20.} Te zwak betekent te ver: de uitdijing is versneld --- donkere
energie, bewijsmatig rustend op de eenvormigheid van de kaarsen (vandaar
de zorg om hun natuurkunde).
\textbf{21.} Het uitsluitingsbeginsel $\to$ \omterm{thm:b3:quantum-statistics:fermi}{ontaardingsdruk} $\to$ een
universele grensmassa $\to$ gestandaardiseerde explosies $\to$ de gemeten
versnelling van het heelal.
\textbf{22.} Volg de wiskunde waarheen zij ook leidt: het ``ongerijmde''
gedrag van sterren voorbij de grens bleek neutronensterren en zwarte gaten
op te leveren --- beide inmiddels bij duizenden gecatalogiseerd.
\textbf{23.} In 1915 stond geen bekende natuurkunde materie een ton per
kubieke centimeter toe; in 1926 maakte de statistiek van Fermi en Dirac
haar tot de standaardtoestand van elk koud dicht gas: de waarneming
wachtte op de statistiek, niet omgekeerd.
\textbf{24.} Dat dwergen over een reeks massa's stabiel bestaan (de tak
$R(M)$), en dat die tak bij $1.4\,M_\odot$ \emph{ophoudt}: beide staan in
dat histogram getekend.
\textbf{25.} $(\hbar c/Gm_{\text{n}}^2)^{3/2}m_{\text{n}} \approx
1.4\,M_\odot$: drie constanten leggen de zwaarste \omterm{ex:b3:quantum-statistics:dwarf}{witte dwerg} vast --- en
daarmee de ontstekingsmassa van de standaardkaarsen die het lot van het
heelal wogen.
\end{solution}
